\subsection{A sufficient pointwise theorem in the original norm}
\label{sec:archimedean-pointwise-assembly}

We now assemble the pointwise estimate at the single quadratic
Gram power needed by the arithmetic consumer. This avoids the
stronger requirement that the estimate hold at every smaller
power in the source's intermediate interface.

Put \(s=r+1\), \(m=2s\), \(n=k+2s-1\), and let
\(\delta=\|\lambda'-\lambda\|_1\). The coefficient \(F\)
has Harish--Chandra parameter \(\lambda'\). Retain the original local
parameter, character, representation and section hypotheses,
including the standard shape, interlacing, separating-gap or
cut alternative, and the noncompact-gap smooth-decay reserve.
Let \(C_J\ge1\) be the constant in the established iterated
leading-minor estimate
\[
 jN^2(D_s^L F)\le b^{2C_Jm^2L}jN^2(F),
 \qquad \delta+s(L+s+1)\le q_0 .                       \tag{PA1}
\]
The existing moment theorem also requires all gaps of the base
parameter to be at least \(C_g^{\rm mom}Lm\). For the powers
used below, \(Lm\le96s^3\), so this follows from the original
\(C_gm^3\) bound after a fixed enlargement of \(C_g\).
Here \(q_0\) is the original minimum of the noncompact and cloud
gap reserves, with both original divisors retained. The estimate
and its actual integer-family construction are given in
\texttt{global-norm-frontier-iterated-minors.tex} and
\texttt{global-norm-frontier-jacquet-gram.tex}.\cite{current-construction-record}

\begin{theorem}[Pointwise bound with an explicit sufficient reserve]
\label{thm:archimedean-pointwise-sufficient}
Under the original local hypotheses, take the all-gap constant
above a fixed threshold depending only on \(C_h\) and the
previously established analytic supplier constants. Assume also
\[
 \operatorname{AllGaps}(\lambda,200\delta),\qquad
                    q_0\ge\delta+s(48s^2+s+1).         \tag{PA2}
\]
For every original finite compact-finite coefficient \(F\), every
admitted radial section, and every represented relevant point
\(\gamma\) with \(G_L(\gamma)>0\), one has
\[
 n_L(\gamma)\operatorname{orbSq}_L(F,\gamma)
 \le b^{C_*m^4}\max\{1,G_L(\gamma)^{-1}\}^{C_*m^2}
                              \mathcal N_L(F).          \tag{PA3}
\]
The same conclusion holds at any nonzero-Gram point, interpreting
the original orbital supremum as zero if there is no representative.
All norms, point weights, characters and quotient measures in
(PA3) are the original ones. The constant is independent of
\(r,k,F,\gamma\), and the section. With the preceding choice
of \(C_J\), the deliberately generous value
\[
               C_*=64+16\max(C_h,0)+24C_J              \tag{PA4}
\]
is sufficient. The actual character is \(\kappa=1\) or
\(\kappa=\mathrm{ka}\), as in the local consumer.
\end{theorem}

\begin{proof}
We first use the standard character and the original normalized
analytic tower. Set
\[
 p=6s^2,\qquad
 a=\begin{cases}0,&u\ne u',\\30s^2+6,&u=u',\end{cases}
 \qquad L=2p+a\le48s^2,
 \qquad d=sL\le48s^3 .                                \tag{PA5}
\]
The equal-type case has the original separating adjacent gap,
so the two compact shifts of length \(a\) exist once its fixed
gap constant is at least \(36\). Their parameters stay between
the original adjacent endpoints, hence retain the bound \(b\).
Each shifted pair has a cut of width \(a\); the unshifted
different-type pair has its prescribed cut. Both contain the
quadratic grid of Section~\ref{actual-adaptive-pointwise}.
At the exceptional index, precisely the two inequalities omitted
in (NC3) may change; all its other interlacing inequalities remain
the original ones. The adjacent gap preserves (NC4).

Choose
\[
                         c=\delta+d+1 .                \tag{PA6}
\]
The original all-gap bound \(C_gm^3\), together with the first
condition in (PA2), implies \(4c+8\le\operatorname{gap}_{\min}
(\lambda)\) after a fixed enlargement of \(C_g\).
Indeed the maximum of \(8C_gs^3\) and \(200\delta\) is at
least \(4C_gs^3+100\delta\), which dominates
\(4\delta+192s^3+12\) for \(C_g\ge51\).
Since every \(|\lambda_j|\le b\), this also gives
\(c+2\le b\). The original shape bound gives \(k\le b\)
and \(b\ge3\). In particular all spectral shifts of degree
at most \(d\) stay in the same regular chamber.

In the standard-character realization the original height
hypotheses give
\(b^{-2C_hm}\le T,T'\le b^{2C_hm}\); replacing \(C_h\)
by its nonnegative part is harmless. Take
\(\nu(t)=\prod_j(|t+\lambda_j|+c)\).
Lemma~\ref{lem:archimedean-normalizer-comparison} supplies
the grid scale \(\Lambda_0=b^{(6+2\max(C_h,0))m}\).
Condition (PA2) supplies the original shift family needed in
Theorem~\ref{thm:archimedean-spectral-word-energy}, and
(PA6) supplies its precise spectral majorant.

The original noncompact-gap hypothesis and smooth-decay theorem
give a common decay exponent large enough for both the fixed
Gram-cutoff exhaustion and the level-zero pairing. To check the
order of the constants, put \(X=3sn^2+n^3\). The previously
proved exhaustion needs at most a fixed multiple of \(s^4X\)
after the polynomial degree loss, whereas (PA5) loses at most
\(48s^3\) powers. The original smooth-decay lower bound is
\((C_g/2)s^4X-1\). A fixed increase of \(C_g\) supplies both
budgets. It is independent of the number of differentiations:
all finite derivative seminorms have that same decay exponent.

Let \(A\) be the identity in the different-type case, or one
of the two cofactor-polynomial maps in the equal-type case.
It has degree at most \(as\). Apply the compact pointwise
theorem first with a spatial exhaustion and the fixed positive
Gram cutoff
\[
                 H_\varepsilon=\chi(G/\varepsilon)G^pFA.
\]
The established fixed-cutoff exhaustion gives the same inequality
for this finite-decay coefficient. At a fixed positive-Gram point,
\(\varepsilon<G(\gamma)\) makes the cutoff identically one
on the entire orbital integration locus. Its point energy is
therefore exactly \(G(\gamma)^{2p}E(FA;\gamma)\).
There is no limiting passage from the big cell to another stratum.

Theorem~\ref{pbw-level-zero-cutoff} applies with word budget
\(N=3s\), since
\[
             5+(4s-2)3s=12s^2-6s+5<2p .
\]
It gives convergence of the actual quadratic level-zero forms.
Theorem~\ref{thm:archimedean-spectral-word-energy} bounds their
limits by the same \(jN^2(G^pFA)\). There are only finitely
many words. Letting \(\varepsilon\) decrease to zero in
(NC15) gives
\[
 E(FA;\gamma)
 \le b^{(8+2\max(C_h,0))m^2}G(\gamma)^{-2p}
                                      jN^2(G^pFA).     \tag{PA7}
\]
No constant used only in proving convergence survives this limit.

If \(u\ne u'\), use \(0\le G\le|D_s|^2\) on the radial
domain and then (PA1) with \(L=2p\).
If \(u=u'\), use the two-shift orbital and norm comparisons
of Section~\ref{actual-wigner-orbital-transport}. The factors
\((T')^a\) cancel exactly, and (PA1) applies with \(L=2p+a\).
In either case, with an unnecessary factor two in the first,
\[
 E(F;\gamma)
 \le 2(3b)^{2a}b^{(8+2\max(C_h,0))m^2+2C_Jm^2L}
                   G(\gamma)^{-L}jN^2(F).              \tag{PA8}
\]
For \(b\ge3\), \(a\le36s^2=9m^2\), and \(m\ge2\),
the factor before \(G^{-L}\) is at most
\[
                   b^{(13+\max(C_h,0)+24C_J)m^4}.
                                                        \tag{PA9}
\]

We record the exact conversion from this analytic tower to the
literal local orbital supremum. If \(h_\gamma\) is the number
of complex stabilizer roots, then \(h_\gamma\le s^2\).
Let \(c_N>0\) be the ratio of the two original unipotent Haar
conventions and put \(A_N=(\kappa_Nc_N)^4>0\).
The same two unipotent integrations occur in the orbital and
radial maps. In each complex stabilizer root, the standard
canonical measure is \(\pi^{-1}\) times Lebesgue measure;
the quotient Weil identity gives the reciprocal factor.
The resulting identities are
\[
 \begin{split}
 n_L(\gamma)\operatorname{orbSq}_L(F,\gamma)
      &=A_N\pi^{2h_\gamma}E(F;\gamma),\\
 \mathcal N_L(F)&=A_NjN^2(F).
 \end{split}                                           \tag{PA10}
\]
Unitary frame identifications preserve Hilbert--Schmidt norms.
Weil uniqueness identifies every admitted quotient measure at a
fixed representative. Compact and unipotent fibre covariance,
together with those unitary identifications, identifies the
values for different representatives. Thus (PA10) includes their
supremum. These are the explicit root and pair-measure
calculations in the original orbital-square normalization and
standard Bessel scale comparison.\cite{current-construction-record}
This calculation neither bounds \(A_N\) nor divides by a norm:
multiply (PA8) by \(A_N\) and use both identities.
Since \(\pi^{2h_\gamma}\le b^{m^4}\), its only loss is one
in the exponent of (PA9). Also \(L\le12m^2\), so
\(G^{-L}\le\max(1,G^{-1})^{12m^2}\).

Finally, for the actual character \(\kappa=\mathrm{ka}\), use
Proposition~\ref{prop:archimedean-character-transport}. Its exact
frame transform preserves the representation parameters and the
reserve (PA2). The original bounds
\(a,a^{-1}\le b^{\max(C_h,0)}\), together with \(L\le48s^2\),
bound the extra factor by
\[
 b^{2\max(C_h,0)rs(1+L)}
                  \le b^{7\max(C_h,0)m^4}.             \tag{PA11}
\]
Every exponent containing \(k\) has already cancelled in that
proposition. The constant (PA4) dominates (PA9), the root factor,
(PA11), and the pole coefficient \(12\).

If a required period map or radial norm is zero, the all-cell
factorization gives zero orbital evaluations; all preceding
homogeneous inequalities still hold. A represented point has
nonnegative Gram by the cofactor identity. Hence at negative
Gram there are no admitted representatives and the orbital
supremum is zero. This proves the asserted nonzero-Gram version.
\end{proof}

The reserve (PA2) is part of this theorem and must be checked in
its application. In particular, an all-gap bound proportional
to \(\delta\) alone has not been substituted for the original
\(q_0\) condition. The ensuing arithmetic application supplies
this reserve for its entire actual finite shift family.

\begin{corollary}[The original unshifted pointwise consumer]
\label{cor:archimedean-unshifted-pointwise}
In the original arithmetic range with nonempty cloud-gap set,
the pointwise estimate (PA3) holds for \(\lambda'=\lambda\)
under the original standard hypotheses, after a fixed increase
of their gap threshold. No additional spectral-reserve premise
is needed in this unshifted case.
\end{corollary}
\begin{proof}
Let \(C_q\ge1\) be the fixed constant in the original definition
of \(q_0\). The original all-gap bound \(C_gm^3\) and
noncompact-gap bound imply
\[
 q_0\ge C_gm^3/C_q^2=8C_gs^3/C_q^2.
\]
The nonempty cloud-gap condition is essential in this step;
it holds in the actual pinned arithmetic range.
Since \(s(48s^2+s+1)\le50s^3\), the fixed choice
\(C_g\ge(25/4)C_q^2\), in addition to the preceding fixed
thresholds, supplies (PA2) at \(\delta=0\).
Apply Theorem~\ref{thm:archimedean-pointwise-sufficient}.
\end{proof}
